\subsection{A rational pendant reduction for a three-object triangle}

\begin{lemma}[Proved: a singleton pendant at an acyclic corner]
\label{CP2:A-rational-pendant-deletion}
Let $p$ be an object of a finite category $C$ whose only
strictly adjacent object is $q$. Suppose their unique nonempty
strict hom-set is a singleton and $BC(q,q)$ is rationally
acyclic. Let $E$ be the full subcategory obtained by deleting
$p$. Then the inclusion $BE\to BC$ is a rational homology
equivalence and induces a rational cohomology-ring isomorphism.
There is no restriction on the monoid $C(p,p)$.
\end{lemma}
\begin{proof}
There are no arrows between $p$ and any object other than $q$,
and arrows between $p$ and $q$ go in one direction. A nerve
simplex that contains $p$ can therefore use no object besides
$p$ and $q$: any further strict object change would force an
additional arrow incident to $p$ by composition. Thus
\[
 BC=BE\cup B C\{p,q\},\qquad
 BE\cap B C\{p,q\}=BC(q,q),
\]
as a union of nerve subcomplexes. The singleton cross hom-set
makes the two-object category the categorical join of its
two corner monoids. Consequently
\[
 B C\{p,q\}\simeq BC(p,p)*BC(q,q).
\]
The join with a nonempty rationally acyclic space is rationally
acyclic, by the reduced homology formula for a join over the
field $\mathbb Q$. Both the intersection and the attached
piece therefore have zero positive rational homology and are
connected. The Mayer--Vietoris sequence proves that
$BE\to BC$ is a rational homology equivalence. The induced
cohomology map is multiplicative, proving the ring assertion.
\end{proof}

\begin{corollary}[Proved: a singleton pendant over a singleton-long triangle]
\label{CP2:A-pendant-singleton-long-triangle}
Suppose a finite four-object category has a singleton pendant
as in the preceding lemma. If its three retained objects
form full linear acyclic support, their long hom-set is a
singleton, and their three corner classifying spaces are
rationally acyclic, every positive rational cohomology product
of the whole category is zero. In particular it suffices
that the retained three corner orders are at most four.
\end{corollary}
\begin{proof}
The pendant's neighbor is retained, so its corner is rationally
acyclic. The lemma identifies the rational cohomology ring
with that of the retained triangle. Theorem
\ref{CP2:A-baseline-singleton-long} annihilates every positive
product in that ring. The corner-order consequence uses
Lemma~\ref{CP2:A-five-types}.
\end{proof}

For example, the frozen four-object eighteen-arrow profile
with index $175$ has hom-set cardinality matrix
\[
 \begin{pmatrix}
 1&0&0&1\\
 0&1&2&1\\
 0&0&3&7\\
 0&0&0&2
 \end{pmatrix}.
\]
Delete the pendant object $0$. The retained triangle has
corner orders $1,3,2$ and singleton long hom-set $1\to3$,
so the corollary excludes the entire profile without any
composition enumeration or any assumption on the sink's
two-element monoid table.
